\par\bigskip
\pdfbookmark[1]{Remarks}{purae-bemerkungen}
\begin{center}\textsc{Remarks.}\end{center}
\begingroup\small
Only a few articles remain of the original continuation of this paper from Article~19 onward, which was devoted to special cases. They deal with the quadratic equation whose roots are the two periods of $\frac{n-1}2$ terms; the manuscript stops at the beginning of the investigation that was to determine the sign of the square root occurring in its solution. The agreement of this surviving beginning with the paper Summatio quarumdam serierum singularium shows that the author changed his plan to make the determination of the sign just mentioned the subject of a separate paper. Comparison with the citation in Article~8, where the manuscript mistakenly gave the third edition of \textsc{Lagrange}'s work instead of the second, shows that this manuscript dates from 1808. That publication of the preceding material had not been abandoned is shown by a leaf evidently inserted later at Article~19, on which a different continuation begins and which already refers to the paper Summatio, etc., for the determination of the sign. This second continuation, which also soon stops, is presented here. In preparing the edition, the text of the manuscript, which was entirely ready for printing, was faithfully retained; only in Article~16 did the formulas for the second case have to be added.
\begin{flushright}R. \textsc{Dedekind}.\end{flushright}
\endgroup
\BookRule{.34\textwidth}
